\section*{الفصل \ref{ch:b2:anthropogenic-climate-change} --- تغيّر المناخ الناجم عن الإنسان}

\begin{solution}{exo:b2:anthropogenic-climate-change:1}
(أ) $3.75 - 0.60 + 0.22 + 0.45 + 0.11 = \qty{3.93}{W.m^{-2}}$، إذ تضيف التكيّفات
\qty{0.18}{W.m^{-2}}. (ب) $3.93/3.71 = 1.059$: أي أكثر بنسبة \qty{5.9}{\percent}. (ج) تكيّف درجة حرارة
التروبوسفير، $-0.60$: فالتروبوسفير الأدفأ تبعث مزيدًا من الأشعّة تحت الحمراء نحو الفضاء، وهذا ينزع
جزءًا من الطاقة الإضافية قبل أن يستجيب السطح.
\end{solution}

\begin{solution}{exo:b2:anthropogenic-climate-change:2}
(أ) $4\pi(\num{6.371e6})^2 = \qty{5.10e14}{m^2}$. (ب) على مدى 47 سنة، $\num{434.9e21}/(47 \times
\num{3.156e7} \times \num{5.10e14}) = \qty{0.575}{W.m^{-2}}$؛ وعلى مدى 12 سنة، $\num{152.4e21}/(12
\times \num{3.156e7} \times \num{5.10e14}) = \qty{0.789}{W.m^{-2}}$، وهما القيمتان 0.57 و0.79 في
التقييم. (ج) $0.91 \times 434.9 = \qty{396}{ZJ}$.
\end{solution}

\begin{solution}{exo:b2:anthropogenic-climate-change:3}
(أ) $-3.22 + 1.30 + 0.35 + 0.42 - 0.01 = \num{-1.16}$، وهي القيمة الصافية. (ب) $3.93/1.16 =
\qty{3.39}{\celsius}$؛ و$3.93/3.22 = \qty{1.22}{\celsius}$. (ج) $3.39/1.22 = 2.78$: فالتغذيات الراجعة
تكاد تضاعف احترار بلانك ثلاث مرّات.
\end{solution}

\begin{solution}{exo:b2:anthropogenic-climate-change:4}
(أ) بيتاغرام واحد من الكربون يساوي $44.009/12.011 = 3.664$~Gt من ثنائي أكسيد الكربون: $0.45 \times 3.664 =
\qty{1.65}{\celsius}$ لكلّ \qty{1000}{PgC}، ويصبح المدى من 0.989 إلى 2.31. (ب) $2390 \times
12.011/44.009 = \qty{652}{PgC}$. (ج) $0.45 \times 2.39 = \qty{1.08}{\celsius}$.
\end{solution}

\begin{solution}{exo:b2:anthropogenic-climate-change:5}
(أ) $g = 159/398 = 0.400$؛ و$T = (398/\num{5.670e-8})^{1/4} = \qty{289}{K}$. (ب) $-0.6 \times 4
\times \num{5.670e-8} \times 288^3 = \qty{-3.25}{W.m^{-2}.K^{-1}}$؛ وعند \qty{298}{K}، $-3.60$.
(ج) أقلّ حساسيةً: فالاستجابة تزداد بحسب $T^3$، فيتخلّص كوكب أدفأ من كلّ درجة إضافية
بسهولة أكبر، إذا تساوت التغذيات الراجعة الأخرى.
\end{solution}

\begin{solution}{exo:b2:anthropogenic-climate-change:6}
(أ) $1.85 - 0.50 = \qty{1.35}{W.m^{-2}.\celsius^{-1}}$. (ب) $\sqrt{0.32^2 + 0.20^2} = 0.38$.
(ج) المجموع معروف أفضل بكثير من كلّ من الحدّين: فخطآهما مرتبطان ارتباطًا عكسيًا قويًا. فحيث
تحترّ التروبوسفير العليا أكثر، تحتفظ أيضًا ببخار أكثر، فتأتي تغذية راجعة سالبة أقوى لمعدّل الانخفاض
الحراري مع تغذية راجعة موجبة أقوى لبخار الماء.
\end{solution}

\begin{solution}{exo:b2:anthropogenic-climate-change:7}
(أ) $\lambda = (0.59 - 2.20)/1.03 = \qty{-1.56}{W.m^{-2}.\celsius^{-1}}$؛ و$3.93/1.56 =
\qty{2.52}{\celsius}$. (ب) $3.93/(1.56 - 0.5) = 3.93/1.06 = \qty{3.71}{\celsius}$. (ج) $3.93/1.46 =
\qty{2.69}{\celsius}$. ويجد التقييم، بالتوزيعات الكاملة، 3.5 و\qty{2.7}{\celsius}.
\end{solution}

\begin{solution}{exo:b2:anthropogenic-climate-change:8}
(أ) $0.617/3682 = \qty{1.68e-4}{\celsius^{-1}}$. (ب) $0.113 \times 0.396 = \qty{0.0447}{m}$، أي
\qty{44.7}{mm}، وهي قريبة من القيمة \qty{47.5}{mm} في الميزانية. (ج) $\num{2e-4} \times 0.2 \times 700 =
\qty{0.028}{m}$، أي \qty{28}{mm}.
\end{solution}

\begin{solution}{exo:b2:anthropogenic-climate-change:9}
(أ) $\sum xy = 0.311$، و$\sum x^2 = 0.300$: $\beta = 1.04$. (ب) مجال \qty{95}{\percent} هو
$1.04 \pm 1.96 \times 0.08$، أي من 0.88 إلى 1.20: وهو يستبعد 0، فالاستجابة مكتشَفة، ويحتوي على
1، فالتغيّر المرصود متوافق مع التغيّر المحاكى. (ج) المجال من 0.34 إلى 0.66 يظلّ
يستبعد 0 (مكتشَفة)، لكنه يستبعد 1: فالنموذج يبالغ في تقدير الاستجابة بعامل قدره نحو
اثنين.
\end{solution}

\begin{solution}{exo:b2:anthropogenic-climate-change:10}
(أ) $\sum g^2 = 1.20$، و$\sum a^2 = 0.18$، و$\sum ga = -0.46$، و$\sum gy = 0.74$، و$\sum ay = -0.28$.
والمعادلتان الطبيعيتان هما $1.20\,\beta_g - 0.46\,\beta_a = 0.74$ و$-0.46\,\beta_g + 0.18\,\beta_a
= -0.28$؛ ومحدّدهما $1.20 \times 0.18 - 0.46^2 = 0.0044$، و$\beta_g = (0.74 \times
0.18 - 0.46 \times 0.28)/0.0044 = 1.00$، و$\beta_a = (-1.20 \times 0.28 + 0.46 \times 0.74)/0.0044 =
1.00$. (ب) غازات الدفيئة $+\qty{0.80}{\celsius}$، والهباء الجوي $-\qty{0.30}{\celsius}$. (ج) المحدّد
صغير لأن $g$ و$a$ متناسبتان تقريبًا (ارتباطهما $-0.99$): فيمكن امتصاص قليل
من الضجيج في $y$ برفع معاملَي القياس معًا، فيكون كلّ منهما ضعيف
التقييد مع أن أثرهما المجتمع معروف جيّدًا.
\end{solution}

\begin{solution}{exo:b2:anthropogenic-climate-change:11}
(أ) $(0.93 - 0.1)/0.45 \times 1000 = \qty{1844}{Gt}$؛ ناقص 210 و$26 \times 0.93 = 24$:
\qty{1610}{Gt}. (ب) $0.93 - x = 0.45 \times (1350 + 210 + 24)/1000 = 0.713$، فيكون $x =
\qty{0.22}{\celsius}$: أي عند الطرف الأعلى للمدى المقيَّم من 0.1 إلى \qty{0.2}{\celsius} تقريبًا. (ج) لا
تتوقّف الميزانية على الاحترار غير الناجم عن CO$_2$، $x$، إلا عبر $-x/\text{TCRE}$: فكلّ
\qty{0.1}{\celsius} يكلّف $0.1/0.45 \times 1000 = \qty{222}{Gt}$، أيًّا كان الحدّ.
\end{solution}

\begin{solution}{exo:b2:anthropogenic-climate-change:12}
(أ) $3.93/1.81 = \qty{2.17}{\celsius}$، و$3.93/0.51 = \qty{7.71}{\celsius}$، و$3.93/1.16 =
\qty{3.39}{\celsius}$. (ب) طول الجانب الأدنى $3.39 - 2.17 = \qty{1.22}{\celsius}$، وطول الجانب الأعلى
$7.71 - 3.39 = \qty{4.32}{\celsius}$: فالدالة $1/|\lambda|$ محدّبة، فالخطوات المتساوية في
$\lambda$ نحو الصفر ترفع الحساسية أكثر فأكثر. (ج) قيمة $\lambda$ سالبة سلبيةً ضعيفة تعطي حساسية
كبيرة جدًّا، والتغيّرات الصغيرة في $\lambda$ قرب الصفر يصعب استبعادها؛ أمّا الحساسية المنخفضة فتحتاج إلى $\lambda$
سالبة سلبيةً قوية، تحدّها \omterm{def:b2:anthropogenic-climate-change:feedbacks}{استجابة بلانك} ويناقضها الاحترار المرصود.
\end{solution}

\begin{solution}{pb:b2:anthropogenic-climate-change:1}
\textbf{1.} $2.16 + 0.54 + 0.21 + 0.41 = \qty{3.32}{W.m^{-2}}$.
\textbf{2.} $-0.22 - 0.84 = \qty{-1.06}{W.m^{-2}}$.
\textbf{3.} $3.32 + 0.47 + 0.05 - 1.06 - 0.20 + 0.08 + 0.06 = \qty{2.72}{W.m^{-2}}$، وهو المجموع
المقيَّم.
\textbf{4.} $3.93\log_2(410/278) = \qty{2.20}{W.m^{-2}}$، في حدود \qty{2}{\percent} من 2.16.
\textbf{5.} $3.93\log_2(427.35/278) = \qty{2.44}{W.m^{-2}}$.
\textbf{6.} حدّ التفاعل بين الهباء الجوي والسحب، من $-1.45$ إلى $-0.25$، أي بعرض \qty{1.20}{W.m^{-2}}، مقابل 1.52 لعرض
المجموع: فهو يحمل معظم شكّ المجموع لأنه كبير وضعيف
التقييد.
\textbf{7.} $0.79 \times \num{5.10e14} \times 12 \times \num{3.156e7} = \qty{1.53e23}{J}$، أي
\qty{153}{ZJ}.
\textbf{8.} $0.91 \times 153 = \qty{139}{ZJ}$.
\textbf{9.} $\lambda = (0.59 - 2.20)/1.03 = \qty{-1.56}{W.m^{-2}.\celsius^{-1}}$.
\textbf{10.} $3.93/1.56 = \qty{2.52}{\celsius}$.
\textbf{11.} $3.93 \times 1.03/2.20 = \qty{1.84}{\celsius}$.
\textbf{12.} $3.93/1.16 = \qty{3.39}{\celsius}$، أي أعلى بمقدار \qty{0.87}{\celsius}.
\textbf{13.} كان الاحترار حتى الآن أشدّ حيث التغذيات الراجعة أشدّ سلبيةً؛ ومع تطوّر
النمط يصبح معامل التغذية الراجعة أقلّ سلبيةً، فتتجاوز الحساسية البعيدة المدى
الحساسيةَ الفعلية.
\textbf{14.} $0.45 \times 44.009/12.011 = \qty{1.65}{\celsius}$ لكلّ \qty{1000}{PgC}.
\textbf{15.} من 0.989 إلى \qty{2.31}{\celsius} لكلّ \qty{1000}{PgC}.
\textbf{16.} $2390 \times 12.011/44.009 = \qty{652}{PgC}$.
\textbf{17.} $0.45 \times 2.39 = \qty{1.08}{\celsius}$.
\textbf{18.} في حدود التقريب، يفسّر ثنائي أكسيد الكربون وحده الاحترار البشري كلّه: فاحترار
غازات الدفيئة الأخرى وتبريد الهباء الجوي يكادان يتعادلان.
\textbf{19.} $1.5 - 1.07 = \qty{0.43}{\celsius}$.
\textbf{20.} $(0.43 - 0.1)/0.45 \times 1000 = \qty{733}{Gt}$ من ثنائي أكسيد الكربون.
\textbf{21.} $733 - 210 = \qty{523}{Gt}$.
\textbf{22.} $26 \times 0.43 = \qty{11.2}{Gt}$؛ و$523 - 11.2 = \qty{512}{Gt}$.
\textbf{23.} تنخفض بمقدار $0.1/0.45 \times 1000 = \qty{222}{Gt}$، إلى \qty{290}{Gt}: وهو المقدار
\qty{220}{Gt} الذي يقول التقييم إن الخيارات بشأن المحرّكات الأخرى تحرّك به الميزانية.
\textbf{24.} محسوبةً من الاستجابة TCRE، تبلغ \omterm{def:b2:anthropogenic-climate-change:budget}{ميزانية الكربون المتبقّية} لحدّ \qty{1.5}{\celsius} ابتداءً من
مطلع 2020 \textbf{\qty[text-series-to-math=true]{512}{Gt}} من ثنائي أكسيد الكربون، في حدود
\qty{3}{\percent} من القيمة المقيّمة \qty{500}{Gt} باحتمال \qty{50}{\percent}.
\end{solution}
